\documentclass[11pt,a4paper]{ctexart}
\usepackage[margin=2.2cm]{geometry}
\usepackage{xcolor}
\usepackage{enumitem}
\usepackage{amssymb}
\usepackage{booktabs}
\usepackage[normalem]{ulem}
\usepackage{array}
\usepackage{longtable}
\usepackage{titlesec}
\usepackage{hyperref}
\definecolor{navy}{HTML}{1A365D}
\hypersetup{colorlinks=true,linkcolor=navy,urlcolor=navy}
\setlist{nosep}
\titleformat{\section}{\Large\bfseries\color{navy}}{}{0pt}{}
\titleformat{\subsection}{\large\bfseries\color{navy}}{}{0pt}{}
\pagestyle{plain}
\renewcommand{\arraystretch}{1.25}

\title{\textbf{九年级下册语法总结（网传版）}\\[4pt]
\large 人教版九年级下册英语（网传电子版）· 按教材书末 Grammar 附录整理 · 仅供教师参考}
\author{初中英语资料站}
\date{2026-11-12 · 附配套专项练习 192 题}
\begin{document}
\maketitle
\begin{center}
\fbox{\parbox{0.95\textwidth}{\small\textbf{⚠ 网传版本 · 非官方教材 · 仅供教师参考}\\
本册教材尚未正式出版（官方版预计 2027 年发布），本讲义依据网传电子版整理，板块划分、表格条目与例句可能存在偏差。\textbf{学生请勿使用}；仅供教师备课参考。正式内容请以 2027 年官方教材为准。}}
\end{center}
\bigskip

\section*{使用说明}
\begin{itemize}[leftmargin=1.4em]
\item 本讲义按九年级下册（网传版）教材书末 Grammar 附录整理，表格与例句取自教材原文。
\item 网页版：\url{https://1901151458-prog.github.io/english-materials/g9x-grammar/}
\item ⚠ 网传版本，非官方教材，学生请勿使用，仅供教师参考。
\end{itemize}

\section{动词的时态（Verb Tenses）}
\noindent\textcolor{gray}{六种主要时态一览}\par\smallskip
动词的时态是用以表示各种时间和动作方面的动词形式。这里主要介绍\textbf{一般现在时、一般过去时、一般将来时、现在进行时、过去进行时和现在完成时}六种动词时态。

{\footnotesize
\begin{longtable}{p{4.47cm}p{4.47cm}p{4.47cm}}
\hline
时态 & 基本用法 & 例句 \\
\hline
\endfirsthead
\hline
时态 & 基本用法 & 例句 \\
\hline
\endhead
\hline
\endfoot
一般现在时 & 表示经常发生的动作或存在的状态 & I often go to school at 7 o’clock.\\ She usually does her homework after school.\\ Earth moves round the sun. \\
一般过去时 & 表示过去某一时间发生的动作或存在的状态 & Both students enjoyed the school trips and chib activities.\\ We went to the cinema last night. \\
一般将来时 & 表示单纯的将来的事实 & I will miss you all.\\ I will become a full-time worker at the factory after graduation. \\
现在进行时 & 表示现在正备进行的动作 & What are they doing about the problem?\\ It’s raining hard. \\
过去进行时 & 表示过去某一时间正在进行的动作 & Last year, I was still studying Chinese in Beijing.\\ I was watching TV when you called me yesterday. \\
现在完成时 & 表示动作发生在过去，但对现在有影响或产生了结果 & The company has already offered me a job.\\ Have you ever seen a Dragon Boat race? \\
\end{longtable}
}

\section{情态动词（Modal Verbs）}
\noindent\textcolor{gray}{表推测与可能性：must / can’t / may·might·could}\par\smallskip
情态动词 \textbf{must、might、may、could、can} 等可以用于表达推测和可能性，但所包含的意义不尽相同。

\subsection*{1. must}

must 表示\textbf{很大的可能性}，意为「一定；必定；肯定」，多用于\textbf{肯定句}中。

\begin{itemize}[leftmargin=1.4em]
\item That movie \underline{must be} very interesting. Shall we go to see it?\\ 那部电影一定很有意思。
\item The cakes look nice. They \underline{must be} delicious.\\ 蛋糕看起来不错，一定很好吃。
\end{itemize}

\subsection*{2. can 和 can’t}

can 表示「\textbf{可能}」，常用于否定句或疑问句中，用于疑问句时带有「惊异；怀疑；不相信」等含义。\textbf{can’t 表示「不可能」}。

\begin{itemize}[leftmargin=1.4em]
\item — \underline{Can} the news be true?\\ — That \underline{can’t} be true.\\ 这消息可能是真的吗？——不可能是真的。
\item It looks so thin. It \underline{can’t} be heavy.\\ 它看起来这么薄，不可能重。
\end{itemize}

\subsection*{3. may、might、could}

may、might、could 都能表示「\textbf{可能}」。\textbf{could、might 语气更委婉，might 语气最为缓和}，表示不太确定的推测或可能性。

\begin{itemize}[leftmargin=1.4em]
\item He \underline{may / might come} to visit us tomorrow.\\ 他明天可能会来看我们。
\item He thinks Beijing \underline{could be} the first place in Asia where humans used fire.\\ 他认为北京可能是亚洲最早使用火的地方。
\item The evidence shows that they \underline{might also have developed} the skill of making silk.\\ 证据表明他们可能也发展了缫丝技术。
\end{itemize}

\begin{center}\fbox{\parbox{0.92\textwidth}{\small \textbf{推测强度：}\textbf{must}（一定）＞ \textbf{may}（可能）＞ \textbf{might / could}（也许，最不确定）；\textbf{can’t} 表示「不可能」。}}\end{center}

\section{冠词（Articles）}
\noindent\textcolor{gray}{定冠词 the、不定冠词 a/an、零冠词}\par\smallskip
冠词是一种虚词，常放在名词前面，帮助说明名词所指的人或事物的含义。冠词本身不能单独使用，在句中一般不重读。冠词可分为\textbf{定冠词、不定冠词和零冠词}三类。

\subsection*{1. 定冠词 the}

不论是单数名词还是复数名词，也不论是可数名词还是不可数名词，其前均可用定冠词 the。

{\footnotesize
\begin{longtable}{p{6.70cm}p{6.70cm}}
\hline
用法 & 例词或例句 \\
\hline
\endfirsthead
\hline
用法 & 例词或例句 \\
\hline
\endhead
\hline
\endfoot
特指某人或某（些）事物 & Let’s go and buy something for the party. \\
指双方都知道的人或事物 & There’s a large sports field next to the school hall. \\
指上文提过的人或事物 & The file is sent to the 3D printer. \\
用在世界上独一无二的事物前 & The sun is bigger than the moon. \\
用在序数词和形容词最高级前 & The blue whale is the largest and the heaviest animal in the world.\\ Da Vinci was one of the first Western artists to study the human body. \\
用在自普通名词构成的专有名词前 & the Louvre Pyramid, the Atlantic Ocean \\
用在姓氏的复数形式前 & the Wright brothers \\
用在一些习惯用语中 & all the same, in the morning / afternoon / evening, on the left / right, the day before yesterday \\
\end{longtable}
}

\subsection*{2. 不定冠词 a / an}

a / an 是不定冠词，\textbf{a 用在发音以辅音开头的单词之前}（a book、a girl），\textbf{an 用在发音以元音开头的单词之前}（an hour、an English teacher）。不定冠词\textbf{只用于可数名词单数前}。

{\footnotesize
\begin{longtable}{p{6.70cm}p{6.70cm}}
\hline
用法 & 例词或例句 \\
\hline
\endfirsthead
\hline
用法 & 例词或例句 \\
\hline
\endhead
\hline
\endfoot
指人或事物的某一种类 & It’s a big farm.\\ I want to be an engineer in the future. \\
指某人或某物，但不具体说明何人何物 & Every year, my family and I take a photo together under the big tree in our village.\\ A true friend is always there for you. \\
表示「一」这个数量，但数的概念没有 one 强烈 & Our town has just built a new hospital and library. \\
用于某些固定词组中 & a few, a little, a lot of, a moment \\
\end{longtable}
}

\subsection*{3. 零冠词}

英语中不少场合下名词前面不需要使用冠词，有的语法书把这种现象称作「零冠词」。

{\footnotesize
\begin{longtable}{p{6.70cm}p{6.70cm}}
\hline
用法 & 例词或例句 \\
\hline
\endfirsthead
\hline
用法 & 例词或例句 \\
\hline
\endhead
\hline
\endfoot
在专有名词和不可数名词前 & Class 1, China, Nobel Prize, Mount Huangshan, water \\
名词前已有作定语的 this、these、my、your、some、any、each、every 等词 & How much is this birthday cake?\\ These ping-pong bats are nice.\\ It comes with some candles.\\ He is proud of his job. \\
复数名词表示一类人或事物时 & Most of the rubbish comes from homes, schools, hospitals, and businesses.\\ The noise from ships and underwater machines can confuse animals like whales and dolphins. \\
在星期、月份、节日前 & Monday, January, May Day, Christmas \\
在称呼语、头衔和职务等名称前 & This is Professor Liu.\\ — Who’s your class teacher? — Ms Gao. \\
在一日三餐前 & For lunch, she likes hamburgers, salad, and oranges. \\
在球类运动的名称前 & Around 1:00, he meets his friends, and they play football. \\
用在某些固定词组中 & at home / school, by train, go to school / bed, on time, at night / noon, after class / school, be late for school \\
\end{longtable}
}

\section{句子类型（Types of Sentences）}
\noindent\textcolor{gray}{简单句 / 并列句 / 复合句}\par\smallskip
句子按其结构可分为\textbf{简单句、并列句和复合句}三种类型。

{\footnotesize
\begin{longtable}{p{4.47cm}p{4.47cm}p{4.47cm}}
\hline
类别 & 构成 & 例句 \\
\hline
\endfirsthead
\hline
类别 & 构成 & 例句 \\
\hline
\endhead
\hline
\endfoot
简单句 & 由一个主语（或并列主语）和一个谓语（或并列谓语）构成的句子 & My mother often goes shopping a weekends.\\ Some students walk or ride bicycles to school. \\
并列句 & 由并列连词（and、so、but、or 等）把两个或两个以上的简单句连在一起而构成的句子 & They are sweet rice balls, and we eat them during the Lantern Festival.\\ School starts at eight, so I get up at six thirty.\\ There are many boatyaces around the world, but this one is special. \\
复合句 & 由一个主句和一个或一个以上的从句构成的句子 & I wonder if the two water festivals are similar.\\ What do you want to be when you grow up?\\ Do you have everything that you need? \\
\end{longtable}
}

\begin{center}\fbox{\parbox{0.92\textwidth}{\small \textbf{快速判断：}只有一个主谓结构 → \textbf{简单句}；两个或多个简单句用并列连词连接 → \textbf{并列句}；主句里套了从句（宾语从句、状语从句、定语从句等）→ \textbf{复合句}。}}\end{center}

\section{构词法（Word Formation）}
\noindent\textcolor{gray}{合成法 / 派生法 / 转化法 / 缩写和简写}\par\smallskip
英语中很多单词的构成形式具有规律，掌握单词的构成规律有助于理解和记忆词汇。常见构词法有\textbf{合成、派生和转化}。缩写和简写也是构词法的一种。

\subsection*{1. 合成法}

由两个或两个以上的词合成一个新词，这种构词法叫作合成法。

{\footnotesize
\begin{longtable}{p{6.70cm}p{6.70cm}}
\hline
类型 & 例 \\
\hline
\endfirsthead
\hline
类型 & 例 \\
\hline
\endhead
\hline
\endfoot
复合名词 & newspaper（名词+名词）、hotline（形容词+名词） \\
复合动词 & highlight（形容词+动词）、overcome（副词+动词）、sunbathe（名词+动词） \\
复合形容词 & hard-working（形容词+动词-ing 形式）、man-made（名词+动词-ed 形式）、worldwide（名词+形容词） \\
复合副词 & anyway（不定代词+名词）、downstairs（副词+名词）、everywhere（形容词+副词） \\
复合代词 & somebody、anything（不定代词+名词） \\
复合数词 & fifty-four、one-third（数词+数词） \\
\end{longtable}
}

\subsection*{2. 派生法}

在一个单词前面或后面加上一个词缀构成新词，这种构词法叫作派生法。\textbf{加在单词前的词缀叫前缀，加在单词后的词缀叫后缀。}

{\footnotesize
\begin{longtable}{p{6.70cm}p{6.70cm}}
\hline
前缀 & 示例 \\
\hline
\endfirsthead
\hline
前缀 & 示例 \\
\hline
\endhead
\hline
\endfoot
un-（不、非，表示反向或反义） & unlucky, unhappy, untidy, unusual \\
dis-（不、反面，表示否定） & disagree, disappoint, dishonest, dislike \\
bi-（双、二、表示数量） & bicycle, bilateral, bilingual, binoculars \\
inter-（相互、交互、在一起） & international, internet, interview \\
re-（再次、重新） & recycle, retell, reuse, rewrite \\
tele-（远距离的） & telephone, telescope, television \\
\end{longtable}
}

{\footnotesize
\begin{longtable}{p{6.70cm}p{6.70cm}}
\hline
后缀 & 示例 \\
\hline
\endfirsthead
\hline
后缀 & 示例 \\
\hline
\endhead
\hline
\endfoot
-er / -or（表示「从事某种职业的人」，名词后缀） & farmer, reporter, researcher, worker, actor, director, visitor \\
-ist（表示「主义者，从事某种职业者」等，名词后缀） & artist, communist, dentist, pianist, scientist, socialist, tourist \\
-ness（表示「处于某种状态」，名词后缀） & happiness, illness, kindness, sadness \\
-ion / -tion（表示「行为；行为的过程或结果、状态」等，名词后缀） & action, attraction, conversation, education, graduation, invention \\
-ful（表示「充满……的；具有……的」，形容词后缀） & beautiful, careful, forgetful, grateful, harmful, peaceful \\
-y（表性质，形容词后缀） & dirty, funny, healthy, noisy, sandy, shiny, smelly, starry, sticky, windy \\
-ing（形容词后缀） & boring, exciting, interesting, moving, outgoing, tiring, touching \\
-ed（形容词后缀） & bored, excited, gifted, interested, skilled, tired, used \\
-al（形容词后缀） & formal, international, natural, normal, personal, traditional \\
-able（表示「能……的；具有……性质的」，形容词后缀） & comfortable, enjoyable, recyclable, unforgettable \\
-less（表示「没有；无」，形容词后缀） & careless, helpless, homeless, weightless \\
-ly（副词或形容词后缀） & greatly, really, widely, daily, friendly, lovely \\
\end{longtable}
}

\subsection*{3. 转化法}

一个单词由一种词类转换为另一种词类，这种构词法叫转化法。单词转化后的意义往往与之前的意义联系密切。

{\footnotesize
\begin{longtable}{p{6.70cm}p{6.70cm}}
\hline
转化类型 & 例 \\
\hline
\endfirsthead
\hline
转化类型 & 例 \\
\hline
\endhead
\hline
\endfoot
动词转化为名词 & cry v. 哭 → (have a) cry n. 哭（了一顿）\\ look v. 看 → (have a) look n. 看（一下） \\
名词转化为动词 & book n. 书 → book (a seat) v. 预定（座位）\\ water n. 水 → water (the flower) v. 浇（花） \\
形容词转化为动词 & slow adj. 慢的 → slow (down) v. 缓（下来）\\ calm adj. 镇静的 → calm (yourself) v. 使（自己）平静 \\
\end{longtable}
}

\subsection*{4. 缩写和简写}

缩写和简写主要采取「截头」「去尾」或者「一截截头又去尾」的方法生成新词。

{\footnotesize
\begin{longtable}{p{3.35cm}p{3.35cm}p{3.35cm}p{3.35cm}}
\hline
缩写/简写 & 原词 & 缩写/简写 & 原词 \\
\hline
\endfirsthead
\hline
缩写/简写 & 原词 & 缩写/简写 & 原词 \\
\hline
\endhead
\hline
\endfoot
phone & telephone & plane & airplane \\
lab & laboratory & ad & advertisement \\
exam & examination & flu & influenza \\
\end{longtable}
}

另外，还有很多缩写词由各个单词的首字母组成：

{\footnotesize
\begin{longtable}{p{3.35cm}p{3.35cm}p{3.35cm}p{3.35cm}}
\hline
缩写词 & 全称 & 缩写词 & 全称 \\
\hline
\endfirsthead
\hline
缩写词 & 全称 & 缩写词 & 全称 \\
\hline
\endhead
\hline
\endfoot
AI & Artificial Intelligence & BCE & before the Common Era \\
CE & Common Era & CPC & Communist Party of China \\
GDP & gross domestic product & PE & physical education \\
PRC & People’s Republic of China & TCM & Traditional Chinese Medicine \\
UN & United Nations &  &  \\
\end{longtable}
}

\section{配套专项练习（网传版）}
\noindent\textcolor{gray}{四个专项各 48 题、共 192 题 ⚠ 学生请勿使用，仅供教师参考}\par\smallskip
\paragraph{动词时态专项 · 48 题}

\paragraph{情态动词与冠词专项 · 48 题}

\paragraph{句子类型专项 · 48 题}

\paragraph{构词法专项 · 48 题}

\begin{center}\fbox{\parbox{0.92\textwidth}{\small \textbf{⚠ 网传版提醒：}四个专项合计 192 题，题目依据网传电子版教材设计，可能存在偏差。\textbf{学生请勿使用}，仅供教师备课参考。\\ \textbf{考前速查：}一页纸速查表 →（A4 单页，可打印）}}\end{center}

\end{document}
